\paragraph{The first global Selmer bound.}
Put $q=|k|$ and
$\nu=\binom a2+\binom{n-a}2+s\,a(n-a)$, where $s$ is the size
of the odd family. Retain the arithmetic tower and its local maps,
the ordinary hypotheses, distinguished rigidity and auxiliary simple
spectrum used above. Its arithmetic system supplies finite global
duality, the local Euler identities and the unramified count. For
Galois groups with finite allowed ramification including the primes
above $\ell$, these are the classical Poitou--Tate and Euler theorems
\cite[I, Theorems 2.3, 2.8, 4.10, 5.1 and Corollary 4.15]{milne-adt}.
Realizing the supplied groups and local maps as that arithmetic
system remains part of the arithmetic construction.
Its chosen real involutions and distinguished involution satisfy
$\chi(c_v)=\chi(c)=-1$, as part of the retained arithmetic data.
At an ordinary place, $d_v$ is its local Euler degree, and
$d_v'+1$ is the rank of the torsion-free pro-$\ell$ abelian quotient
used by the original diagonal condition. For the actual local
Galois group, $d_v'=d_v=[F_v:\mathbf Q_\ell]$.

\begin{proposition}[First Selmer estimate]
\label{prop:selmer-first-global-bound}
Let the arithmetic system have $d$ real places and $q_{\rm aux}$
auxiliary places, with ordinary degrees satisfying
$\sum_vd_v=\sum_vd_v'=d$. There is a constant $C$ such that, for
every $m\geq0$ and all sufficiently large $N$, the original first
Selmer cocycle set is finite and
\[
 \#\operatorname{Sel}^{1}_{N,m}
 \leq \#B^{1,\mathrm{pol}}_{N,m}\,
               q^{mn(d+q_{\rm aux})+C}.
\]
Here $B^{1,\mathrm{pol}}_{N,m}$ is its group of boundaries with skew
primitives. One may take
\[
          C=C_G+C_{\rm det}+C_0+\sum_vC_v+\sum_jC_j,
\]
using the invariant, detecting and ordinary constants already proved.
\end{proposition}
\begin{proof}
The case $m=0$ is immediate. For $m\geq1$, the primal and paired
modules are Tate dual under
\[
 (X,Y)\longmapsto\operatorname{Tr}_{k/\mathbb F_\ell}
                       [t^{m-1}\eta]\operatorname{tr}(XY).
\]
Matrix units, nilpotent generators and powers of $t$ pair in complementary
coordinates. The finite-field trace pairing is perfect; its value at
1 need not be nonzero. Conjugation and transpose prove equivariance.

Normalizing a global cocycle to be zero at $c$ identifies
$\operatorname{Sel}^{1}_{N,m}/B^{1,\mathrm{pol}}_{N,m}$ with the
cohomological Selmer group for $M_m$. Its ordinary condition is the
image of $Z_{v,m}$, its old unramified condition is unramified
cohomology, and its distinguished and auxiliary conditions are
unrestricted. The normalizing boundary preserves these conditions.

Write $L_v$ for the local cohomology conditions. The finite global
duality inequality bounds this Selmer group by
\[
 \#H^0(\Gamma_N,M_m)\,
 \#H^1_{\rm rel}(\Gamma_N,W_m)
 \prod_{v\ \mathrm{finite}}\frac{\#L_v}{\#H^0(G_v,M_m)}
 \prod_{v\mid\infty}\frac1{\#M_m^{c_v}}.
\]
The relaxed dual group retains only the boundary conditions at the
unrestricted primal places. This inequality follows from the usual
Selmer-dual Selmer cardinality identity by discarding the denominator
$\#H^0(\Gamma_N,W_m)$ and the other dual local conditions.

At an ordinary place, let $\mathfrak b_{v,m}$ be the actual
transported upper triangular primitive module. It has
$q^{m(n+\nu)}$ elements, and its boundaries belong to $Z_{v,m}$.
The preceding filtration counts its $n$ diagonal homomorphisms
using $d_v'+1$ and its $\nu$ strictly upper coordinates using
the local Euler degree $d_v$.
Writing $B_v$ for all local boundaries gives
\[
 \frac{\#L_v}{\#H^0(G_v,M_m)}
 =\frac{\#(B_v+Z_{v,m})}{\#M_m}
 \leq\frac{\#Z_{v,m}}{\#\mathfrak b_{v,m}}
 \leq q^{m(nd_v'+\nu d_v)+C_v}.
\]
For the middle inequality, the boundary map on $\mathfrak b_{v,m}$
has kernel $\mathfrak b_{v,m}\cap H^0(G_v,M_m)$; insert its image
in $B_v\cap Z_{v,m}$. Unramified factors are one, since a finite
module has equally many Frobenius invariants and coinvariants.
At the distinguished and auxiliary places, local Euler degree zero
and local duality identify the factor with $\#H^0(G_v,W_m)$.
Proposition~\ref{prop:selmer-three-invariant-bounds} and eventual
matching bound these by $q^{C_0}$ and $q^{mn+C_j}$, respectively.

For a real involution $c_v$, put $P_v=R(c_vc)$. The original
scalar conjugate self-duality and $\chi(c_vc)=1$, which follows
from the two retained signs, give $P_v^T=P_v$. The action is
$Z\mapsto-P_vZ^TP_v^{-1}$. Multiplication by $P_v$ identifies
its fixed module with the skew part of $M_m$, which has $\nu$
free coordinates. Both $P_v$ and its inverse preserve the primal
coefficient algebra. Thus $\#M_m^{c_v}=q^{m\nu}$ exactly.

The global invariant factor is at most $q^{C_G}$. The relaxed dual
group injects into the auxiliary-local quotient bounded by
Proposition~\ref{prop:selmer-tower-detection}, so its order is at most
$q^{C_{\rm det}}$. Multiplication of the displayed factors leaves
the exponent
\[
 m\left(n\sum_vd_v'+\nu\sum_vd_v+nq_{\rm aux}-\nu d\right)+C
                         =mn(d+q_{\rm aux})+C.
\]
The arithmetic finiteness theorem and the finite boundary group give
finiteness of the original cocycle set. All constants were chosen
before $m$; only the eventual level threshold depends on it
\cite{alpo-first-global-selmer}.
\end{proof}

This finishes the first Selmer estimate under its retained arithmetic
inputs. The second estimate and the construction of the patched
modules are separate steps.
